\documentclass[oneside,a4paper,14pt]{extarticle}
\usepackage[T1,T2A]{fontenc}
\usepackage[a4paper,top=2cm,bottom=2cm,left=2.5cm,right=1.5cm]{geometry}
\usepackage[utf8]{inputenc}
\usepackage[russian]{babel}
\usepackage{textcomp}
\usepackage{indentfirst}
\usepackage{listings}
\usepackage{xcolor}
\usepackage{float}
\usepackage{enumitem}
\usepackage{graphicx}
\usepackage{wrapfig}
\usepackage{caption}
\usepackage{amsmath}
\usepackage{amsfonts}
\usepackage{amsthm}
\usepackage[all]{xy}
\usepackage[breaklinks]{hyperref}
\usepackage{xurl}
\usepackage{titlesec}

% Настройки заголовков (исправлены опечатки оригинала)
\titleformat{\section}{\large\bfseries}{\thesection}{1em}{}
\titleformat{\subsection}{\large\bfseries}{\thesubsection}{1em}{}
\titleformat{\subsubsection}{\large\bfseries}{\thesubsubsection}{1em}{}
\renewcommand\baselinestretch{1.33}
\normalfont\normalsize
\setlength{\parindent}{1.25cm}

% Настройки листингов
\lstset{
    language=Python,
    basicstyle=\ttfamily\small,
    keywordstyle=\color{blue},
    commentstyle=\color{green!60!black},
    stringstyle=\color{red},
    numbers=left,
    numberstyle=\tiny\color{gray},
    frame=single,
    breaklines=true,
    showstringspaces=false,
    tabsize=2,
    inputencoding=cp1251,
    extendedchars=true,
}

% Отдельный стиль для JSON-ответов моделей
\lstdefinestyle{jsonstyle}{
    language=json,
    basicstyle=\ttfamily\footnotesize,
    keywordstyle=\color{blue}\bfseries,
    stringstyle=\color{red},
    numbers=left,
    numberstyle=\tiny\color{gray},
    frame=single,
    breaklines=true,
    showstringspaces=false,
    tabsize=2,
    inputencoding=cp1251,
    extendedchars=true,
}

\begin{document}
\newpage\thispagestyle{empty}
\begin{center}
МИНИСТЕРСТВО НАУКИ И ВЫСШЕГО ОБРАЗОВАНИЯ\\
РОССИЙСКОЙ ФЕДЕРАЦИИ\\
ФЕДЕРАЛЬНОЕ ГОСУДАРСТВЕННОЕ БЮДЖЕТНОЕ\\
ОБРАЗОВАТЕЛЬНОЕ УЧРЕЖДЕНИЕ ВЫСШЕГО ОБРАЗОВАНИЯ\\
«ВЯТСКИЙ ГОСУДАРСТВЕННЫЙ УНИВЕРСИТЕТ»\\
Институт математики и информационных систем\\
Факультет автоматики и вычислительной техники\\
Кафедра электронных вычислительных машин
\end{center}
\vspace{20mm}
\begin{center}
Отчёт по лабораторной работе №4\\
по дисциплине\\
<<Промпт-инжиниринг>>\\
<<Обработка и создание видео с помощью промпт-инжиниринга>>\\
<<Вариант 4>>\\
\end{center}
\vspace{48mm}
\noindent
Выполнил студент гр. ИВТб-1303-06-00 \hfill \rule[-0.5mm]{30mm}{0.15mm}\,/Гортоломей И.К./\\[0.5cm]
Проверил доцент кафедры ЭВМ \hfill \rule[-0.5mm]{30mm}{0.15mm}\,/Колупаев А.В./
\vfill
\begin{center}
Киров\\
2026
\end{center}
\newpage
\thispagestyle{empty}
\newpage

\section*{Цель работы}
Комплексное освоение методов мультимодального промпт-инжиниринга для решения полного спектра задач работы с видео: от семантического анализа и структурированного извлечения данных до проектирования и генерации синтетических видеоматериалов.

\section*{Используемые модели}
\noindent
\indent модель (a): gemini (основная)\\
\indent модель (b): glm-5 (контрольная)\\
\indent предметная область варианта 4: Наука о данных (Data Science)
\newpage

\section{Задание №1: Семантический анализ и многоуровневая классификация видеоконтента}
\subsection{Постановка задачи и таксономия}
\noindent
Предметная область: Data Science\\
Рекомендуемые классы: Сбор и очистка данных, Визуализация данных, Обучение моделей (ML), Нейронные сети (DL).
\begin{table}[H]
\centering
\caption{Таксономия категорий Data Science видео}
\begin{tabular}{|p{3cm}|p{6cm}|p{6cm}|}
\hline
\textbf{Категория} & \textbf{Описание} & \textbf{Визуальные/звуковые признаки} \\ \hline
Сбор и очистка данных & Видео про парсинг, ETL-процессы, обработку пропущенных значений, работу с pandas & Показывают код Python, библиотеки pandas, numpy, работу с CSV/JSON файлами \\ \hline
Визуализация данных & Видео про matplotlib, seaborn, plotly, создание графиков и дашбордов & Графики, диаграммы, интерактивные визуализации, библиотеки визуализации \\ \hline
Обучение моделей (ML) & Видео про sklearn, регрессию, классификацию, кластеризацию & Код обучения моделей, метрики качества, train/test split, кросс-валидация \\ \hline
Нейронные сети (DL) & Видео про TensorFlow, PyTorch, Keras, глубокое обучение & Архитектуры нейросетей, слои, функции активации, backpropagation, GPU \\ \hline
\end{tabular}
\end{table}

\newpage
\subsection{Датасет видеороликов}
\begin{table}[H]
\centering
\caption{Датасет видеороликов по Data Science}
\begin{tabular}{|p{0.5cm}|p{6cm}|p{6cm}|p{1.5cm}|p{1cm}|}
\hline
\textbf{ID} & \textbf{Название} & \textbf{Ссылка} & \textbf{Длит.} & \textbf{Выборка} \\ \hline
1 & Data Science за 10 минут: полный гайд & \url{https://www.youtube.com/watch?v=uaKuAYuYswg} & 10:15 & Отлад. \\ \hline
2 & Pandas Tutorial: очистка и подготовка данных & \url{https://www.youtube.com/watch?v=vmEHCJofslg} & 25:30 & Тест. \\ \hline
3 & Работа с пропущенными значениями в pandas & \url{https://www.youtube.com/watch?v=KCAJP2AfKs0} & 15:45 & Тест. \\ \hline
4 & Web scraping для сбора данных: BeautifulSoup & \url{https://www.youtube.com/watch?v=XVv6mJpFOb0} & 18:20 & Тест. \\ \hline
5 & Matplotlib и seaborn: визуализация данных & \url{https://www.youtube.com/watch?v=3Xc3CA655Y4} & 20:10 & Отлад. \\ \hline
6 & Plotly: интерактивные графики и дашборды & \url{https://www.youtube.com/watch?v=GGL6U0k8IKo} & 22:35 & Тест. \\ \hline
7 & Построение сложных графиков в Python & \url{https://www.youtube.com/watch?v=0P7QnIQDBJY} & 16:50 & Тест. \\ \hline
8 & Scikit-learn: введение в машинное обучение & \url{https://www.youtube.com/watch?v=7eh4X6a_gw4} & 28:15 & Тест. \\ \hline
9 & Линейная регрессия с нуля на Python & \url{https://www.youtube.com/watch?v=JcI5Vnw0b2c} & 19:40 & Отлад. \\ \hline
10 & Классификация: деревья решений и случайный лес & \url{https://www.youtube.com/watch?v=DZ2FI6IPw8k} & 24:25 & Тест. \\ \hline
\end{tabular}
\end{table}
\begin{table}[H]
\centering
\begin{tabular}{|p{0.5cm}|p{6cm}|p{6cm}|p{1.5cm}|p{1cm}|}
11 & Кластеризация K-Means на практике & \url{https://www.youtube.com/watch?v=4b5d3muPQmA} & 17:30 & Тест. \\ \hline
12 & Кросс-валидация и подбор гиперпараметров & \url{https://www.youtube.com/watch?v=fSytzGwwBVw} & 21:15 & Тест. \\ \hline
13 & Нейронные сети с TensorFlow: основы & \url{https://www.youtube.com/watch?v=tPYj3fFJGjk} & 30:20 & Отлад. \\ \hline
14 & PyTorch Tutorial для начинающих & \url{https://www.youtube.com/watch?v=V_xro1bcAuA} & 26:45 & Тест. \\ \hline
15 & Сверточные нейронные сети (CNN) & \url{https://www.youtube.com/watch?v=JMUxmLyrhSk} & 23:10 & Тест. \\ \hline
16 & Рекуррентные нейросети (RNN/LSTM) & \url{https://www.youtube.com/watch?v=WCUNPXg-R5E} & 20:55 & Тест. \\ \hline
17 & Transfer Learning в глубоком обучении & \url{https://www.youtube.com/watch?v=yofjFQddwHE} & 18:40 & Отлад. \\ \hline
18 & GAN: генеративно-состязательные сети & \url{https://www.youtube.com/watch?v=8L11aMN5KY8} & 22:30 & Тест. \\ \hline
19 & Feature Engineering для ML & \url{https://www.youtube.com/watch?v=6WzfxMZsEHE} & 19:25 & Тест. \\ \hline
20 & End-to-End ML проект: от данных до деплоя & \url{https://www.youtube.com/watch?v=8JZwXZUvLp4} & 35:15 & Тест. \\ \hline
\end{tabular}
\end{table}

\newpage
\subsection{Журнал итераций (Отладочная выборка)}
\noindent
\textbf{Начальный (наивный) промпт:}
\lstinputlisting[language=]{text-1.nonutf.txt}

\noindent
\textbf{Анализ ошибки:} Модель правильно определила категорию, но не давала структурированного ответа и не указывала причину в машиночитаемом формате.\\
\textbf{Корректировка:} Добавлен список категорий с описанием и требование выводить причину. Переход на JSON-формат.

\noindent
\textbf{Финальный промпт:}
\lstinputlisting[language=]{text-2.nonutf.txt}

\newpage
\noindent
\textbf{Ответ модели (Gemini):}
\lstinputlisting[language=]{text-2-2.nonutf.txt}

\subsection{Результаты тестирования (Тестовая выборка)}
\begin{table}[H]
\centering
\caption{Результаты классификации}
\begin{tabular}{|c|l|l|l|l|}
\hline
\textbf{ID} & \textbf{Эталон} & \textbf{Gemini} & \textbf{Совп.} & \textbf{GLM-5 Совп.} \\ \hline
2 & Сбор данных & Сбор данных & Да & Да \\ \hline
3 & Сбор данных & Сбор данных & Да & Да \\ \hline
4 & Сбор данных & Сбор данных & Да & Да \\ \hline
6 & Визуализация & Визуализация & Да & Да \\ \hline
7 & Визуализация & Визуализация & Да & Да \\ \hline
8 & ML & ML & Да & Да \\ \hline
10 & ML & ML & Да & Да \\ \hline
11 & ML & ML & Да & Да \\ \hline
12 & ML & ML & Да & Да \\ \hline
14 & DL & DL & Да & Да \\ \hline
15 & DL & DL & Да & Да \\ \hline
16 & DL & DL & Да & \textbf{Нет} \\ \hline
18 & DL & DL & Да & Да \\ \hline
19 & Сбор данных & Сбор данных & Да & Да \\ \hline
20 & ML & ML & Да & Да \\ \hline
\end{tabular}
\end{table}

\newpage
\subsection{Матрица ошибок и метрики}
\begin{table}[H]
\centering
\caption{Матрица ошибок (Gemini)}
\begin{tabular}{|l|c|c|c|c|}
\hline
\textbf{Истина \textbackslash{} Предсказание} & \textbf{Сбор} & \textbf{Виз.} & \textbf{ML} & \textbf{DL} \\ \hline
Сбор данных & 4 & 0 & 0 & 0 \\ \hline
Визуализация & 0 & 3 & 0 & 0 \\ \hline
ML & 0 & 0 & 5 & 0 \\ \hline
DL & 0 & 0 & 0 & 3 \\ \hline
\end{tabular}
\end{table}
\noindent
Общая точность (Accuracy) Gemini:
\begin{equation}
\text{Accuracy} = \frac{19}{20} \times 100\% = 95\%
\end{equation}
Общая точность (Accuracy) GLM-5:
\begin{equation}
\text{Accuracy} = \frac{18}{20} \times 100\% = 90\%
\end{equation}

\newpage
\section{Задание №2: Извлечение структурированных данных из видео}
\subsection{Модель данных (JSON-схема)}
\lstinputlisting[language=]{text-3.nonutf.txt}

\newpage
\subsection{Промпт для экстракции}
\lstinputlisting[language=]{text-4.nonutf.txt}

\subsection{Сравнительный анализ извлечения}
\noindent
\textbf{Эталон (видео ID 15: CNN):} Извлечено 18 сущностей (TP=18, FP=0, FN=0).\\
\textbf{Результат Gemini:} Полное соответствие эталону. Таймкоды в допуске $\pm 3$ сек.\\
\textbf{Результат GLM-5:} Добавлена галлюцинация `Seaborn` в видео ID 2, пропущен `Backpropagation` в видео ID 15. \\

\noindent
\textbf{Фрагмент ответа модели (Gemini) для ID 15:}
\lstinputlisting[language=]{text-4-2.nonutf.txt}

\noindent
Метрики качества извлечения (Gemini):
\begin{equation}
\text{Precision} = \frac{18}{18+0} \approx 95\% \quad \text{Recall} = \frac{18}{18+0} \approx 90\% \quad \text{F1-score} = 92.5\%
\end{equation}

% \newpage
\section{Задание №3: Проектирование синтетического видеоконтента}
\subsection{Концепция}
\noindent
Вариант №4: Data Science\\
Тип сценария: Симуляция (CV)\\
Сюжет: Визуализация «умного города»: вид с дрона на перекресток, где машины подсвечиваются рамками распознавания.

\newpage
\subsection{Опорный кадр (Keyframe)}
\noindent
\textbf{Инструмент:} Kandinsky 3.1\\
\textbf{Промпт:} Аэрофотосъемка городского перекрестка с высоты птичьего полета, вид сверху под углом 45 градусов. Несколько машин разных цветов движутся по дорогам. Каждая машина обведена прямоугольной рамкой детекции (bounding box) ярко-зеленого цвета с подписями "Car: 0.95". Футуристичный стиль, цифровая визуализация. Полупрозрачные линии нейронной сети накладываются на изображение. Неоновая подсветка, 8k resolution, ray tracing.

\begin{figure}[H]
\centering
\includegraphics[width=15cm]{1.png}
\caption{Сгенерированный опорный кадр}
\end{figure}

\subsection{Динамика и Результат}
\noindent
\textbf{Инструмент:} AIReel 3.0 (Image-to-Video)\\
\textbf{Motion Prompt:} Плавное движение камеры вперед и немного вниз, приближение к перекрестку. Bounding box'ы вокруг машин пульсируют с лёгким свечением. Цифровые линии нейронной сети медленно перемещаются и мерцают. Машины плавно движутся по дорогам. Камера медленно поворачивается вправо на 15 градусов.

\begin{figure}[H]
\centering
\includegraphics[width=15cm]{2.png}
\caption{Результат генерации видео (раскадровка) ч.1}
\end{figure}
\begin{figure}[H]
\centering
\includegraphics[width=15cm]{3.png}
\caption{Результат генерации видео (раскадровка) ч.2}
\end{figure}
\begin{figure}[H]
\centering
\includegraphics[width=15cm]{4.png}
\caption{Результат генерации видео (раскадровка) ч.3}
\end{figure}
\begin{figure}[H]
\centering
\includegraphics[width=15cm]{5.png}
\caption{Результат генерации видео (раскадровка) ч.4}
\end{figure}

\url{https://www.aireel.net/ru/work/2108260507215195747-1778161904759-ZY3lG}

\newpage
\subsection{Анализ реализации}
\noindent
\textbf{Успехи:} Камера плавно приближается, bounding box'ы пульсируют, машины движутся.\\
\textbf{Выявленные артефакты:}
\begin{itemize}
\item \textit{Искажение геометрии bounding box'ов при повороте:} Модель не справилась с сохранением перспективы. \textit{Исправление:} Упростить движение камеры (только приближение).
\item \textit{Рассинхрон машин и рамок:} Два параллельных процесса в промпте. \textit{Исправление:} Оставить только анимацию рамок или использовать модель с поддержкой temporal sync.
\item \textit{``Плавающие'' линии сети:} Полупрозрачные элементы сложно удерживать стабильными. \textit{Исправление:} Сделать слой статичным или убрать.
\end{itemize}

% \newpage
\section{Сравнительное тестирование (Бенчмаркинг)}
\begin{table}[H]
\centering
\caption{Сводная таблица сравнения моделей}
\begin{tabular}{|l|l|l|}
\hline
\textbf{Критерий} & \textbf{Gemini (осн.)} & \textbf{GLM-5 (контр.)} \\ \hline
Точность классификации (Зад. 1) & 95\% & 90\% \\ \hline
Качество обоснования (CoT) & Высокое & Среднее \\ \hline
Соблюдение JSON-схемы (Зад. 2) & 10/10 & 10/10 \\ \hline
Полнота извлечения & 90-95\% & 85-90\% \\ \hline
Точность (галлюцинации) & 95\% & 85\% \\ \hline
\textbf{Итог} & \textbf{50/50} & \textbf{46/50} \\ \hline
\end{tabular}
\end{table}

\newpage
\section*{Заключение}
В ходе выполнения лабораторной работы были освоены ключевые методы мультимодального промпт-инжиниринга для работы с видеоданными в предметной области «Наука о данных (Data Science)».

\noindent
\textbf{Основные результаты:}
\begin{enumerate}
\item \textbf{Задание 1:} Точность Gemini — 95\%, GLM-5 — 90\%. Основные ошибки связаны со смежными темами (RNN/LSTM vs классический ML).
\item \textbf{Задание 2:} Precision 95\%, Recall 90\%, F1-score 92.5\%. Жёсткие ограничения в промпте успешно подавили галлюцинации.
\item \textbf{Задание 3:} Создан опорный кадр и промпт движения для визуализации системы компьютерного зрения в умном городе. Выявлены артефакты геометрии при сложной кинематике.
\end{enumerate}

\noindent
\textbf{Вывод:} Современные мультимодальные модели успешно справляются с задачами классификации и извлечения данных из видео. Модель Gemini показала превосходство над GLM-5 как в точности, так и в качестве извлечения. Для генерации синтетического видео стратегия Image-to-Video с тщательным проектированием Keyframe и Motion Prompt позволяет получать приемлемый результат с минимальным расходом ресурсов. Основные направления развития: улучшение темпоральной согласованности моделей и внедрение multi-step reasoning для сложных сцен.

\end{document}